\chapter[1951 --- Theiler]{1951: Max Theiler}
\label{ch:1951_maxtheiler}

\section*{Main Contribution}

\textit{Developed a vaccine against yellow fever, saving countless lives in tropical regions and demonstrating that a single vaccination could provide lifelong immunity.}\cite{nobel-1951,lecture-1951}

\section{Official Citation}

for his discoveries concerning yellow fever and how to combat it

\section{Background and Historical Context}

Max Theiler was born on January 30, 1899, in Pretoria, South Africa, into a family with deep roots in the Cape Colony. His father, Sir Godfrey Martin Theiler, was a physician and bacteriologist who had made important contributions to the study of animal diseases in South Africa. His grandfather, Martin Theiler, had been a Swiss-born missionary who emigrated to South Africa in the nineteenth century. Growing up in this scientifically oriented family, Max developed an early interest in medicine and the natural sciences.

Theiler attended the University of Cape Town, where he studied medicine and graduated with honors in 1922. He then moved to England, where he obtained a position at the Hospital for Tropical Diseases in London, and subsequently enrolled at St.\ Thomas's Hospital Medical School, where he received his degree in 1924. His training in London exposed him to the tropical diseases that would become the focus of his career, and he developed a particular interest in yellow fever---a devastating viral disease that was endemic throughout much of tropical Africa and South America.

In 1922, Theiler secured a position at the Rockefeller Institute for Medical Research (now Rockefeller University) in New York, where he would spend the next 33 years conducting research on tropical diseases. The Rockefeller Institute, which had been founded in 1901 by John D.\ Rockefeller, was one of the world's leading research institutions, and its emphasis on basic and applied medical research provided Theiler with an ideal environment for the research that would lead to his Nobel Prize.

The scientific context of the early twentieth century was one of intense interest in the viral diseases that plagued tropical and subtropical regions. Yellow fever, malaria, dengue, and other vector-borne diseases were major causes of morbidity and mortality throughout the world, and the development of effective vaccines and treatments was a public health priority of the highest order. The Rockefeller Institute had a strong tradition of research on tropical diseases, and Theiler's work on yellow fever was part of a broader effort to understand and control these devastating illnesses.

\section{The Discovery/Contribution}

Max Theiler received the Nobel Prize in Physiology or Medicine ``for his discoveries concerning yellow fever and how to combat it.'' His work encompassed the characterization of yellow fever virus strains, the development of the first effective yellow fever vaccine, and the elucidation of the pathogenesis of the disease.

\subsection{The Epidemiology of Yellow Fever}

Yellow fever is a viral disease transmitted by mosquitoes of the genus \textit{Aedes}, particularly \textit{Aedes aegypti}. The disease was one of the great scourges of tropical regions, causing periodic epidemics with high mortality rates. In the Americas, yellow fever had played a significant role in the colonization and settlement of the New World: the disease decimated European settlers in the Caribbean and South America during the sixteenth and seventeenth centuries and contributed to the failure of several colonial ventures.

In the early twentieth century, yellow fever remained a major public health threat. The construction of the Panama Canal (1904--1914) was complicated by yellow fever epidemics among the construction workers, and the disease continued to cause outbreaks in Africa and South America throughout the first half of the twentieth century. The development of an effective vaccine was a public health priority of the highest order.

\subsection{Theiler's Studies on Yellow Fever Virus}

Theiler's research on yellow fever began in the 1920s, when he set out to characterize the virus responsible for the disease. Working at the Rockefeller Institute, he established laboratory protocols for the study of yellow fever virus, including the use of mice as experimental models. He demonstrated that yellow fever virus exists in two principal forms: an African strain (the Asibi strain, isolated from a patient in Uganda in 1927) and a South American strain (the French neurotropic strain, isolated in West Africa and used as the basis for an early vaccine developed by the French virologist Maurice Baudouin).

Theiler showed that these two strains differ in their virulence and in their ability to cause disease in experimental animals. The Asibi strain was highly virulent and caused severe liver damage and hemorrhagic fever in primates, while the French neurotropic strain was less virulent but had a tendency to cause encephalitis (inflammation of the brain) when administered by certain routes. Understanding these differences was crucial for the development of a safe and effective vaccine.

\subsection{The Development of the Yellow Fever Vaccine}

Theiler's primary contribution was the development of the 17D strain of yellow fever vaccine, which remains one of the most effective and widely used vaccines in the world. Theiler attenuated (weakened) the highly virulent Asibi strain by passing it through mouse and chicken embryo tissue cultures. After more than 100 serial passages through these tissues, the virus was sufficiently weakened that it no longer caused disease in humans, but it retained its ability to stimulate a protective immune response.

The 17D vaccine was first tested in human volunteers in 1936, and the results were remarkable. The vaccine produced a strong immune response, with more than 95\% of vaccinated individuals developing protective antibodies within 10 days. The vaccine was safe, with minimal side effects, and it provided long-lasting immunity---estimates suggest that a single dose of the 17D vaccine provides protection for at least 30 to 40 years, and possibly for life.

The 17D vaccine was first used on a large scale in Brazil in 1937, during a yellow fever epidemic, and it was subsequently deployed throughout tropical regions of Africa and South America. The WHO estimates that the 17D vaccine has been administered to more than 600 million people since its introduction, and it has played a central role in the control of yellow fever worldwide.

\subsection{Understanding the Pathogenesis of Yellow Fever}

Theiler's research also contributed significantly to our understanding of the pathogenesis of yellow fever. He demonstrated that the virus preferentially infects liver cells (hepatocytes), causing extensive damage to the liver that leads to the characteristic jaundice (yellowing of the skin) from which the disease derives its name. He also showed that the virus damages the kidneys and can cause hemorrhagic fever---a severe complication characterized by bleeding from the mucous membranes and internal organs.

Understanding the pathogenesis of yellow fever was important not only for the development of the vaccine but also for the clinical management of the disease. The recognition that yellow fever primarily affects the liver guided the development of supportive treatments for patients with the disease, including fluid replacement, blood transfusions, and careful monitoring of liver function.

\section{Impact on Modern Medicine}

Theiler's discovery of the 17D yellow fever vaccine has had a significant impact on global public health, and his research on yellow fever pathogenesis has informed our understanding of other viral hemorrhagic fevers.

\subsection{Yellow Fever Vaccination Programs}

The 17D vaccine remains the cornerstone of yellow fever control efforts worldwide. The WHO recommends routine vaccination for all children living in endemic areas, as well as vaccination for travelers to regions where yellow fever is prevalent. In 2013, the WHO launched the ``Eliminate Yellow Fever Epidemics'' (EYE) strategy, which aims to protect 1 billion people against yellow fever by 2026 through mass vaccination campaigns and strengthened surveillance.

The success of yellow fever vaccination programs has been one of the great public health achievements of the twentieth century. In many countries, yellow fever has been virtually eliminated as a public health threat, and the disease is now largely confined to forested areas of Africa and South America where the virus circulates between mosquitoes and non-human primates.

\subsection{Insights into Viral Hemorrhagic Fevers}

Theiler's research on the pathogenesis of yellow fever has also informed our understanding of other viral hemorrhagic fevers, including dengue, Zika, chikungunya, and Ebola. These diseases, which share many features with yellow fever, are caused by related viruses that damage blood vessels and the liver, leading to bleeding and organ failure. The insights gained from Theiler's work on yellow fever have been essential for the development of diagnostic tests, treatments, and vaccines for these emerging infectious diseases.

\subsection{Vaccine Development Strategies}

Theiler's development of the 17D vaccine established principles for vaccine development that remain relevant today. His approach---attenuating a virulent virus by serial passage through non-human tissues---was one of the first successful applications of the live-attenuated vaccine strategy, which has since been used to develop vaccines for other diseases, including polio (the Sabin vaccine), measles, and rubella. The 17D vaccine also demonstrated that a single vaccine strain could provide protection against a wide range of viral genotypes, a principle that has informed the development of broadly protective vaccines for other viruses.

\subsection{Current Challenges}

Despite the success of the 17D vaccine, yellow fever remains a public health concern in parts of Africa and South America. Climate change and urbanization are expanding the range of the \textit{Aedes aegypti} mosquito, potentially bringing yellow fever to new regions. There have been recent outbreaks of yellow fever in areas that had not seen the disease for decades, highlighting the importance of maintaining high vaccination coverage. Additionally, global vaccine supply constraints have sometimes limited the ability to respond to outbreaks, leading to renewed efforts to increase vaccine production capacity.

\section{Legacy}

Max Theiler's legacy is that of a scientist who devoted his career to the fight against one of humanity's most devastating infectious diseases. His development of the 17D yellow fever vaccine has saved millions of lives and remains one of the major achievements in the history of vaccinology.

Theiler continued his research at the Rockefeller Institute until 1964, when he retired as head of the Division of Tropical Medicine. He received numerous honors during his career, including the Royal Medal of the Royal Society and the Gairdner Foundation International Award. He died on August 11, 1972, at the age of 73.

The story of Max Theiler and the yellow fever vaccine is a testament to the power of basic research to produce practical benefits for humanity. Theiler's systematic approach to understanding the virus, combined with his ingenuity in developing the attenuation technique that led to the 17D vaccine, exemplifies the best traditions of biomedical research. As we face new challenges in infectious disease---from the emergence of novel viruses to the impact of climate change on vector-borne diseases---the principles and approaches that Theiler pioneered continue to guide and inspire researchers around the world.


\section{Modern Developments}

Theiler's yellow fever vaccine (17D) remains one of the most successful vaccines ever developed, providing lifelong immunity with a single dose. Modern yellow fever outbreaks (Brazil 2017-2018) demonstrate the continued importance of vaccination. Understanding of flavivirus biology (yellow fever, dengue, Zika, West Nile) has enabled development of new vaccines and antivirals. The Dengvaxia dengue vaccine, though controversial, represents progress toward controlling this major global health threat.


\section{Bibliography}

\begin{thebibliography}{9}
\bibitem{nobel-1951}
Nobel Prize Outreach, ``The Nobel Prize in Physiology or Medicine 1951,''\emph{NobelPrize.org}.\\
\url{https://www.nobelprize.org/prizes/medicine/1951/summary/}

\bibitem{lecture-1951}
Max Theiler, ``Nobel Lecture,''\emph{NobelPrize.org}. Winner-authored primary source; downloadable from the official Nobel Prize archive.\\
\url{https://www.nobelprize.org/prizes/medicine/1951/theiler/lecture/}
\end{thebibliography}
